\section{Challenge 2: Số Fibonacci modulo}
\label{sec:challenge2}

\begin{requirementbox}{Yêu cầu bài toán}
Với $N$ truy vấn $A_i$ và một modulo cố định $Q$, chương trình cần trả về danh sách $F(A_i) \bmod Q$ theo đúng thứ tự xuất hiện trong input. Ràng buộc lớn của $N$ và $A_i$ khiến việc tính lặp tuyến tính từng số Fibonacci là không khả thi.
\end{requirementbox}

\subsection{Hiện thực đã xây dựng}
Trong \codepath{challenge/challenge2.py}, mỗi truy vấn Fibonacci được xử lý độc lập. Chương trình sử dụng hai kỹ thuật chính. Kỹ thuật thứ nhất là fast doubling, cho phép tính $F(n)$ với độ phức tạp \bigO{\log n}. Kỹ thuật thứ hai là bảng window, trong đó chỉ số $n$ được tách thành phần thấp và phần cao để tái sử dụng các giá trị Fibonacci đã tiền xử lý.

Kết quả được ghi vào \func{multiprocessing.RawArray}. Cách tổ chức này giúp các worker ghi trực tiếp vào vùng kết quả đã được định vị bằng offset, từ đó bảo toàn thứ tự output mà không cần truyền một danh sách lớn về tiến trình chính.

\subsection{Phân tích song song hóa}
Các truy vấn $F(A_i) \bmod Q$ không phụ thuộc lẫn nhau, do đó đơn vị song song tự nhiên là từng đoạn truy vấn. Mỗi worker nhận một đoạn byte chứa nhiều truy vấn, tách token, tính kết quả và ghi vào vùng output tương ứng. Trước khi worker bắt đầu, các cấu trúc dùng chung như modulo và bảng window phải được khởi tạo hoàn chỉnh.

Việc chia theo byte-range giúp giảm chi phí tạo danh sách truy vấn trung gian. Để bảo toàn thứ tự, mỗi chunk đi kèm với \func{out\_start}, là vị trí bắt đầu ghi kết quả trong \func{RawArray}.

\subsection{Thiết kế thuật toán song song}
\begin{lstlisting}[caption={Pseudocode song song cho Challenge 2}]
MAIN(path):
    raw = read_all_bytes(path)
    N, modulo, query_start = parse_header(raw)
    prepare_fibonacci_tables_if_beneficial()

    output = RawArray(size=N)
    chunks = split_query_region(raw, query_start)

    for each chunk in parallel:
        tokens = raw[chunk.start:chunk.end].split()
        for j, token in enumerate(tokens):
            n = int(token)
            output[chunk.out_start + j] = fib_mod(n, modulo)

    return output[:]
\end{lstlisting}

\subsection{Lưu đồ thuật toán}
\begin{figure}[H]
\centering
\begin{tikzpicture}[node distance=0.8cm and 1.1cm, scale=0.9, transform shape]
\node[flowstart] (start) {Bắt đầu};
\node[flowprocess, below=of start] (parse) {Đọc input và phân tích $N$, $Q$};
\node[flowprocess, below=of parse] (table) {Khởi tạo bảng window hoặc chọn fast doubling};
\node[flowprocess, below=of table] (chunks) {Chia vùng truy vấn thành các chunk};
\node[flowprocess, below=of chunks] (workers) {Worker tính $F(A_i) \bmod Q$ cho từng token};
\node[flowmerge, below=of workers] (write) {Ghi kết quả vào \func{RawArray} theo offset};
\node[flowstart, below=of write] (end) {Trả về danh sách kết quả};
\draw[flowarrow] (start) -- (parse);
\draw[flowarrow] (parse) -- (table);
\draw[flowarrow] (table) -- (chunks);
\draw[flowarrow] (chunks) -- (workers);
\draw[flowarrow] (workers) -- (write);
\draw[flowarrow] (write) -- (end);
\end{tikzpicture}
\caption{Luồng xử lý song song của Challenge 2}
\label{fig:challenge2-flow}
\end{figure}

\subsection{Giải thích vai trò các bước}
Tiến trình chính thực hiện các bước có tính toàn cục: đọc $N$ và modulo $Q$, lựa chọn cách tính Fibonacci, khởi tạo bảng window nếu số lượng truy vấn đủ lớn và cấp phát \func{RawArray} cho output. Các bước này phải hoàn tất trước khi tạo worker vì mọi worker đều phụ thuộc vào cùng giá trị modulo và cùng cấu trúc tra cứu.

Sau đó, vùng chứa các truy vấn được chia thành nhiều chunk. Mỗi chunk đi kèm vị trí bắt đầu trong output, nhờ đó worker có thể ghi kết quả của truy vấn thứ $j$ trong chunk vào đúng ô \func{out\_start + j}. Cách ghi này thể hiện tương tác giữa các tiến trình: worker không cần gửi lại toàn bộ danh sách kết quả, cũng không cần khóa đồng bộ, vì mỗi worker được cấp một vùng ghi không giao nhau.

\subsection{Tiềm năng tăng tốc và ví dụ minh họa}
Bài toán có tiềm năng tăng tốc tốt khi $N$ lớn, vì các truy vấn Fibonacci độc lập. Nếu mỗi truy vấn được tính bằng fast doubling với chi phí \bigO{\log A_i}, tổng khối lượng tính toán đủ lớn có thể được chia cho nhiều worker. Bảng window còn giúp giảm chi phí lặp lại khi nhiều truy vấn có chỉ số lớn và cùng modulo.

Ví dụ, với $N=6$ truy vấn và 3 worker, có thể chia thành ba chunk: worker 1 xử lý $A_0,A_1$, worker 2 xử lý $A_2,A_3$, worker 3 xử lý $A_4,A_5$. Nếu input là $3,4,5,6,7,8$ với $Q=10^9$, các worker lần lượt ghi các kết quả $2,3$, $5,8$, $13,21$ vào các vùng đã định sẵn. Khi tiến trình chính đọc lại \func{RawArray}, thứ tự output vẫn là thứ tự ban đầu của input.

\subsection{Dữ liệu, phụ thuộc và chi phí}
\begin{itemize}
  \item \textbf{Phân phối dữ liệu}: input byte và bảng Fibonacci là dữ liệu đọc chung; output là mảng chia sẻ ghi theo vùng không giao nhau.
  \item \textbf{Phụ thuộc dữ liệu}: các truy vấn độc lập; phụ thuộc duy nhất là bước tiền xử lý bảng phải hoàn tất trước khi chạy worker.
  \item \textbf{Cân bằng tải}: nếu độ dài truy vấn tương đối đồng đều, chia theo byte-range tạo tải gần cân bằng.
  \item \textbf{Chi phí giao tiếp}: thấp vì worker ghi trực tiếp vào \func{RawArray}; chi phí chính nằm ở parse token, tính modulo và khởi tạo process.
\end{itemize}

\subsection{Đánh giá hiệu năng}
\benchmarktable{Challenge 2}

\paragraph{Nhận xét.}
\placeholder{Phân tích lợi ích của bảng window so với fast doubling thuần. Cần nêu rõ khi số truy vấn nhỏ, overhead dựng bảng và tạo tiến trình có thể lớn hơn lợi ích song song hóa.}
